\begin{exercisechunk}
\item \textbf{From H\"older approximation to regression risk.}
\label[exercise]{ex:l15-holder-bias-variance}
\exerciseprofile{2}{2}{2}{\excalculation\quad\exinterpretation}{%
Derive the smoothness-dependent regression rate by balancing approximation
and estimation errors.}
Suppose an estimator based on an $m$-dimensional approximation space satisfies
\begin{equation}
 \sup_{f\in\cH_d^\alpha(L)}
 \E\bigl[\|\widehat f_{n,m}-f\|_{L^2(P_X)}^2\bigr]
 \leq A m^{-2\alpha/d}+B\frac mn,
 \label{eq:exercise-holder-risk-bound}
\end{equation}
where $A,B>0$ do not depend on $m$ or $n$.
\begin{enumerate}[(a),beginpenalty=10000]
\item Balance the two terms in
\cref{eq:exercise-holder-risk-bound}. Show that a choice
\[
 m\asymp n^{d/(2\alpha+d)}
\]
gives risk of order $n^{-2\alpha/(2\alpha+d)}$. Explain why rounding $m$ to
a positive integer does not change the rate.
\item Suppose that the target has the form $f(x)=g(W^\top x)$, the matrix
$W\in\R^{d\times s}$ is known, and the same approximation and estimation
bounds hold in the reduced $s$-dimensional coordinates. Derive the choices
\[
 m\asymp n^{s/(2\alpha+s)},
 \qquad
 \E\bigl[\|\widehat f_{n,m}-f\|_{L^2(P_X)}^2\bigr]
 \lesssim n^{-2\alpha/(2\alpha+s)}.
\]
\item Explain why the calculation in part (b) proves neither the
$s$-dimensional rate for the full $d$-dimensional H\"older ball nor the
existence of an algorithm that finds an unknown subspace from noisy
regression data. Your explanation must identify the target-class restriction
and the missing estimation task separately.
\end{enumerate}
\end{exercisechunk}
